\subsection{表示的半单性}

引理 1. 设 \(\mathfrak{g}\) 是半单李代数。\(\mathfrak{g}\) 的伴随表示是半单的。\(\mathfrak{g}\) 的每个理想和每个商代数都是半单的。

设 \(\mathfrak{a}\) 是 \(\mathfrak{g}\) 的一个理想。相对于 Killing 形式，\(\mathfrak{b}\) 中 \(\mathfrak{a}\) 的正交补 \(\mathfrak{g}\) 是 \(\mathfrak{g}\) 的一个理想，而 \(\mathfrak{a} \cap \mathfrak{b}\) 是交换理想（§ 3，第 6 号，命题 7），因而为零。因此 \(\mathfrak{b}\) 是 \(\mathfrak{a}\) 在 \(\mathfrak{g}\) 中的补。又由于 \(\mathfrak{g}\) 的 Killing 形式非退化，它在 \(\mathfrak{a}\) 和 \(\mathfrak{b}\) 上的限制也都非退化（《代数》IX 章§ 4，第 1 号，命题 1 的推论），从而 \(\mathfrak{a}\) 和 \(\mathfrak{b}\) 都是半单的（第 1 号，定理 1 以及§ 3，第 6 号，命题 9）。

引理 2. 设 \(\mathfrak{g}\) 是李代数。以下两个条件等价：

\begin{enumerate}
\def\labelenumi{(\alph{enumi})}
\item
  \(\mathfrak{g}\) 的所有有限维线性表示都是半单的。
\item
  给定 \(\rho\) 的 \(\mathfrak{g}\) 在有限维向量空间 V 上的线性表示，以及余维为 1 的向量子空间 W，满足关系 \(\rho(x)(\mathbf{V}) \subset \mathbf{W}\) 对所有 \(x \in \mathfrak{g}\) 成立；存在 W 的一条补直线，它在 \(\rho(\mathfrak{g})\) 下稳定（因而被 \(\rho(\mathfrak{g})\) 消去）。
\end{enumerate}

显然 (a) 蕴含 (b)。设 (b) 成立。令 \(\sigma\) 是 \(\mathfrak{g}\) 在向量空间 M 上的有限维表示，N 是在 \(\sigma(\mathfrak{g})\) 下稳定的一个向量子空间。令 \(\mu\) 为 \(\mathfrak{g}\) 在 \(\mathcal{L}(\mathrm{M})\) 上由 \(\sigma\) 规范导出的表示（§ 3，第 3 号）；回忆 \(\mu(x) = \mathrm{ad}_{\mathcal{L}(\mathrm{M})}\sigma(x)\)。令 V（分别为 W）为 \(\mathcal{L}(\mathrm{M})\) 中由 M 映入 N 且在 N 上的限制为数乘变换（分别为零）的线性映射组成的子空间；于是 W 在 V 中余维为 1，且 \(\mu(x)(\mathrm{V})\subset\mathrm{W}\) 对所有 \(x\in\mathfrak{g}\) 成立。由条件 (b)，存在 \(u\in\mathrm{V}\)，它被 \(\mu(x)\) 消去（对所有 \(x\in\mathfrak{g}\)），且它在 N 上的限制是非零数乘变换。将 \(u\) 乘以适当的标量后，可以假设 \(u\) 是 M 到 N 的投影算子。说 \(\mu(x).u=0\)，意味着 \(u\) 与 \(\sigma(x)\) 可交换。因此 \(u\) 的核是 N 在 M 中的补空间，并且在 \(\sigma(x)\) 下稳定（对所有 \(x\in\mathfrak{g}\)）。因此 \(\sigma\) 是半单的。

引理 3. 设 \(\rho\) 是 \(g\) 在有限维向量空间 \(V\) 上的线性表示，且 \(W\) 是 \(V\) 的余维为 1 子空间，并且 \(\rho(x)(V) \subset W\) 对所有 \(x \in g\) 成立。则存在 \(W\) 的一条补直线，它在 \(\rho(g)\) 下稳定。

对所有 \(x \in \mathfrak{g}\)，令 \(\sigma(x)\) 为 \(\rho(x)\) 在 W 上的限制。先设 \(\sigma\) 是单的。若 \(\sigma = 0\)，则 \(\rho(x)\rho(y) = 0\) 对所有 \(x,y\)（属于 \(\mathfrak{g}\)）成立，从而 \(\rho(\mathfrak{g}) = \rho(\mathcal{D}\mathfrak{g}) = \{0\}\)，我们的断言显然成立。若 \(\sigma \neq 0\)，令 \(\mathfrak{n}\) 为 \(\sigma\) 的核，令 \(\mathfrak{m}\) 为 \(\mathfrak{n}\) 在 \(\mathfrak{g}\) 中的一个补理想（引理 1）；则 \(\mathfrak{m} \neq \{0\}\)，且 \(\sigma\) 在 \(\mathfrak{m}\) 上的限制是忠实的；在 \(\mathfrak{m}\) 上，\(\sigma\) 的相关双线性形式的限制是非退化的（命题 1），因而可以构造 Casimir 元 \(c\)，它与 \(\mathfrak{m}\) 和 \(\sigma\) 相关。根据§3，第 7 号的命题 12，\(\sigma(c)\) 是 W 的一个自同构。另一方面，\(\rho(c)(\mathrm{V}) \subset \mathrm{W}\)。因此 \(\rho(c)\) 的核 Z 是 W 的一条补直线；由于 \(c\) 属于 \(\mathfrak{g}\) 的包络代数的中心，\(\rho(c)\) 与所有 \(\rho(x)\) 对所有 \(x \in \mathfrak{g}\) 可交换，从而 Z 在 \(\rho(\mathfrak{g})\) 下稳定。

一般情形下，我们对 V 的维数作归纳。令 T 为 W 的一个极小非零稳定子空间。令 \(\rho'\) 为定义在 \(V' = V/T\) 上的商表示。则对所有 \(x \in g\)，关系 \(\rho'(x)(V') \subset W'\) 成立，其中 \(W' = W/T\) 在 \(V'\) 中余维为 1。由归纳假设，存在一条 \(Z'\)，它是 \(W'\) 的补直线并且在 \(\rho'(g)\) 下稳定。它在 V 中的逆像 Z 在 \(\rho(g)\) 下稳定，包含 T，且 \(Z \cap W = T\)；因而 \(\rho(x)(Z) \subset T\) 对所有 \(x \in g\) 成立。由上面已经证明的结果，存在 T 在 Z 中的一条补直线，它在 \(\rho(g)\) 下稳定；这条直线是 W 在 V 中的补直线，证明完毕。

定理 2 (H. Weyl). 半单李代数的每个有限维线性表示都是完全可约的。

这由引理 2 和 3 立即得到。

定义 2. 若李代数 g 的唯一理想是 \{0\} 和 g，且进一步是非交换的，则称 g 为单李代数。

单李代数是半单的。代数 \(\{0\}\) 不是单的。

命题 2. 李代数 g 为半单的充要条件是它是单代数的乘积。

充分性由（第 1 号，备注 3）得到。反之，设 g 是半单的。由于 g 的伴随表示是半单的，g 是极小非零理想 \(a_{1}, \ldots, a_{m}\) 的直和。于是 g 可与这些 \(a_{i}\) 的乘积代数等同（§ 1，第 1 号）。\(a_{i}\) 的每个理想也都是 g 的理想，因而或者为零，或者等于 \(a_{i}\)。另一方面，\(a_{i}\) 是非交换的。因此 \(a_{i}\) 是单李代数。

推论 1. 半单李代数是其单理想 \(\mathfrak{g}_{\mathrm{i}}\) 的乘积。\(\mathfrak{g}\) 的每个理想都是其中某些 \(\mathfrak{g}_{\mathrm{i}}\) 的乘积。

\(g = a_{1} \times \cdots \times a_{m}\)，其中 \(a_{i}\) 是单的。由于 \(a_{i}\) 的中心为零，\(a_{i}\) 在 g 中的中心化子是 \(a_{j}\)（满足 \(j \neq i\)）的乘积。现在令 a 为 g 的一个理想。若它不包含 \(a_{i}\)，则 \(a \cap a_{i} = \{0\}\)，从而 \([a, a_{i}] = \{0\}\)，并且 a 包含于 \(a_{j}\)（满足 \(j \neq i\)）的乘积中。由此 a 是某些 \(a_{i}\) 的乘积。因此，g 的单理想恰好是 \(a_{i}\)。

半单李代数的单理想称为 g 的单分量。

推论 2. 设 \(\mathfrak{g},\mathfrak{g}'\) 是两个李代数，\(\mathfrak{r}\) 和 \(\mathfrak{r}'\) 分别是它们的根基，\(f\) 是从 \(\mathfrak{g}\) 到 \(\mathfrak{g}'\) 上的一个同态。则 \(\mathfrak{r}' = f(\mathfrak{r})\)。

由于 \(f(\mathfrak{r})\) 是可解的，\(f(\mathfrak{r}) \subset \mathfrak{r}'\)。另一方面，\(\mathfrak{g} / \mathfrak{r}\) 是半单的（§ 5，第 2 号，命题 3），所以 \(\mathfrak{g}' / f(\mathfrak{r})\) 是 \(\mathfrak{g} / \mathfrak{r}\) 的一个商的同构代数，因而是半单的（引理 1），从而 \(f(\mathfrak{r}) \supset \mathfrak{r}'\)（§ 5，第 2 号，命题 3）。

备注。 (1) 定理 2 有逆命题：如果 g 的每个有限维表示都是半单的，那么 g 是半单的。因为伴随表示是半单的，所以 g 的每个理想都有补理想，因而可以看作 g 的一个商。若 g 不是半单的，则 g 有一个非零交换商，因而有一个一维商。现在，一维李代数 K 有非半单表示，例如

\[
\lambda \mapsto \left( \begin{array}{c c} 0 & 0 \\ \lambda & 0 \end{array} \right).
\]

\begin{enumerate}
\def\labelenumi{(\arabic{enumi})}
\setcounter{enumi}{1}
\tightlist
\item
  设 g 是 K 上的李代数，σ 是 g 在向量空间 M 上的一个表示。另一方面，设 \(f\) 是把 \(g\) 映入 \(M\) 的 K-线性映射，满足：
\end{enumerate}

\[
f ([ x, y ]) = \sigma (x). f (y) - \sigma (y). f (x)\tag{1}
\]

对所有 \(x, y\) 属于 \(\mathfrak{g}\)。根据§ 1，第 8 号，例 2，给定 \(\sigma\) 和 \(f\) 等价于给定一个同态 \(x \mapsto (f(x), \sigma(x))\)，它把 \(\mathfrak{g}\) 映入 \(\mathfrak{af}(\mathbf{M})\)。另一方面，我们在前引处已看到，\((f(x), \sigma(x))\) 中的元素 \(\mathfrak{af}(\mathbf{M})\) 可规范地与 \(\rho(x)\) 中的元素 \(\mathfrak{gl}(\mathbf{N})\) 等同（其中 \(\mathrm{N} = \mathrm{M} \times \mathrm{K}\)），后者诱导的 \(\sigma(x)\) 作用于 \(\mathrm{M}\)，并把 \(\mathrm{N}\) 的元素 (0, 1) 映到 \(f(x)\)。于是 \(\rho\) 是 \(\mathfrak{g}\) 在 \(\mathrm{N}\) 上的表示，满足 \(\rho(x)(\mathrm{N}) \subset \mathrm{M}\) 对所有 \(x \in \mathfrak{g}\) 成立。

于是，若 g 是半单的，则存在（引理 3）一条 Z，它是 M 在 N 中的补直线，并被 \(\rho(\mathfrak{g})\) 消去。换句话说，存在一个元素 \(m_{0} \in M\)，使得 \((-m_{0}, 1) \in N\) 被 \(\rho(x)\) 消去，也就是说满足

\[
f (x) = \sigma (x). m _ {0}\tag{2}
\]

对所有 \(x \in g\)。

*假设 K = R。设 G 是具有李代数 g 的连通李群。考虑从 G 到 M 的仿射群 A 的一个解析同态 \(\phi\)，其对应的同态 \(x \mapsto (f(x), \sigma(x))\) 取值于 \(\mathfrak{af}(M)\)。上述结果可以解释为：若 g 是半单的，则 \(\phi(G)\) 保持 M 中的一个点不动。令 H 为 GL(N) 中使 N 内所有平行于 M 的线性簇保持稳定的元素所成的集合。存在（§ 1，第 8 号，例 2）从 A 到 H 的一个规范同构 \(\psi\)。令 Z 为 M 在 N 中的一条补直线。说 \(\rho(\mathfrak{g})\) 消去 Z，等价于说 \((\psi \circ \phi)(G)\) 固定 Z 中的各点，从而（考虑到 \(\psi\) 的定义）\(\phi(G)\) 固定 Z 与 M × \{1\} 的交点在 M 上的投影。

